\chapter{MDP Formulation of Gate Assignment}
\label{chap:mdp}

\section{Problem Description}

An airport has $G$ gates available for aircraft turnaround operations and $R$ runways for arrivals.
Flights arrive according to a stochastic schedule over a finite time horizon $T$. Each flight lands on a specific runway and requires assignment to a gate for a duration that depends on both the aircraft type and the taxiing distance between the runway and the assigned gate.
The objective is to minimize waiting times and queue overflows while maximizing throughput.

\section{State Space}

The state at time $t$ is a tuple $s_t = (t, \mathbf{g}_t, \mathbf{q}_t)$ where:

\begin{itemize}
	\item $t \in \{0, 1, \ldots, T\}$: current time index.
	\item $\mathbf{g}_t \in \{0, 1, \ldots, S_{\max}\}^G$: gate occupancy vector.
	\begin{itemize}
		\item $g_i = 0$: gate $i$ is available.
		\item $g_i = s > 0$: gate $i$ occupied, $s$ time steps remaining.
	\end{itemize}
	\item $\mathbf{q}_t$: queue of waiting flights.
	\begin{itemize}
		\item Each flight $f$ is a tuple $(T_f^{\text{arr}}, r_f, \text{type}_f, p_f)$, where:
		\begin{itemize}
			\item $T_f^{\text{arr}}$: Actual arrival time.
			\item $r_f \in \{1, \ldots, R\}$: The runway the flight landed on.
			\item $\text{type}_f$: Aircraft type (determining base service time).
			\item $p_f$: Priority class.
		\end{itemize}
	\end{itemize}
\end{itemize}

\section{Action Space}

At each time step, the decision-maker selects an action $a_t \in \mathcal{A}_t$ where:
%
\begin{equation}
	\mathcal{A}_t = \{(f, g) : f \in \mathbf{q}_t, g \in \text{Available}(\mathbf{g}_t), \text{Compatible}(f, g)\} \cup \{\text{wait}\}
\end{equation}

An action $(f, g)$ assigns waiting flight $f$ to available gate $g$. The "wait" action defers assignment to the next time step.

\section{Transition Dynamics}

The transition from $s_t$ to $s_{t+1}$ involves:

\subsection{Deterministic Time Evolution}

\begin{equation}
	t \to t+1
\end{equation}

\subsection{Gate Occupancy Updates (Heterogeneous Service)}

Gate occupancy is determined by the aircraft's base turnaround time plus a penalty for taxiing distance.
Let $\Delta(r, g)$ be a deterministic distance matrix representing the taxiing time from runway $r$ to gate $g$.

For each gate $i$:
%
\begin{equation}
	g_i^{t+1} = \begin{cases}
		s-1 & \text{if } g_i^t = s > 1 \\
		0 & \text{if } g_i^t = 1 \text{ (service completes)} \\
		\text{BaseService}(\text{type}_f) + \Delta(r_f, i) & \text{if action assigns flight } f \text{ to gate } i
	\end{cases}
\end{equation}
%
where $\text{BaseService}(\text{type}_f)$ is a random variable sampled from the distribution specific to the aircraft type.

\subsection{Stochastic Schedule-Based Arrivals}

Arrivals are modeled as a perturbed schedule rather than a pure Poisson process.

\noindent Let $\mathcal{T} = \{T_1^{\text{sched}}, T_2^{\text{sched}}, \ldots, T_N^{\text{sched}}\}$ be the master schedule of flight arrival times. The actual arrival time $A_k$ for the $k$-th flight is:
%
\begin{equation}
	A_k = T_k^{\text{sched}} + \delta_k
\end{equation}
%
where $\delta_k$ is a stochastic noise term (e.g., $\delta_k \sim \mathcal{N}(0, \sigma^2)$ or an Erlang delay distribution).

\subsection{Queue Updates}

\begin{equation}
	\mathbf{q}_{t+1} = (\mathbf{q}_t \setminus \{f_{\text{assigned}}\}) \cup \{k : A_k = t+1\}
\end{equation}

\section{Reward Function}

The immediate reward penalizes waiting and encourages efficient assignment:
%
\begin{equation}
	R(s_t, a_t) = -c_{\text{wait}} \cdot |\mathbf{q}_t| - c_{\text{overflow}} \cdot \mathbb{I}[|\mathbf{q}_t| > Q_{\max}] + c_{\text{assign}} \cdot \mathbb{I}[\text{assignment made}]
\end{equation}

\section{Assumptions}

\begin{assumption}
	Base service times are independent random variables with finite support, dependent on aircraft type.
\end{assumption}

\begin{assumption}
	The runway $r_f$ for each flight is predetermined by external air traffic control and is observed upon arrival.
\end{assumption}

\begin{assumption}
	Gates are heterogeneous with respect to accessibility. The total gate occupation time includes a taxiing penalty $\Delta(r, g)$ which makes certain gate-runway combinations strictly more costly than others.
\end{assumption}

\section{Complexity Analysis}

The state space size is:
%
\begin{equation}
	|\mathcal{S}| \approx T \cdot (S_{\max} + \max(\Delta))^G \cdot \binom{F_{\max} + Q_{\max}}{Q_{\max}}
\end{equation}
%
The introduction of the distance matrix $\Delta$ increases the effective support of the service time distribution, but does not increase the dimensionality of the state space, preserving the feasibility of the ADP approach.